\documentclass[11pt,a4paper]{article}
\usepackage[utf8]{inputenc}
\usepackage[margin=1in]{geometry}
\usepackage{tcolorbox}
\usepackage{fancyhdr}
\usepackage{hyperref}
\usepackage{xcolor}

\hypersetup{
    colorlinks=true,
    linkcolor=blue!70!black,
    urlcolor=blue!70!black
}

\pagestyle{fancy}
\fancyhf{}
\rhead{AI Lab Experiments Output Record}
\lhead{Roll No: 25CSEAIML046}
\rfoot{Page \thepage}

\title{\textbf{\LARGE Artificial Intelligence Laboratory (AI Lab)}\\
\Large Compilation of All Experiment Execution Outputs}
\author{\textbf{Soumya Ranjan} \\ Roll / Reg No: \texttt{25CSEAIML046} \\ Department of Computer Science \& Engineering (AI \& ML)}
\date{\today}

\begin{document}
\maketitle
\thispagestyle{fancy}
\hrule
\vspace{1em}

\tableofcontents
\newpage

\section{Experiment 1: Breadth First Search (BFS)}
\textbf{Source File:} \texttt{BFS.py} \\[4pt]
\textbf{Objective:} To implement Breadth-First Search (BFS) graph traversal using a queue data structure to visit vertices level by level. \\[4pt]
\textbf{Input Parameters:} 4 edges: (A B), (A C), (B D), (C E); Starting node: A \\[6pt]
\begin{tcolorbox}[colback=gray!5!white,colframe=blue!70!black,title={\textbf{Program Output: Breadth First Search (BFS)}}]
\begin{verbatim}
---Build Your Graph---
How many edges (connection) does your graph has ?Enter each edge separated by a spacce(e.g., A B):
Edge 1Edge 2Edge 3Edge 4Enter the starting node for BFS: 
Your Graph Dictionary: {'A': ['B', 'C'], 'B': ['A', 'D'], 'C': ['A', 'E'], 'D': ['B'], 'E': ['C']}
Staring BFS Traversal...
Exploring node1: A
Exploring node1: B
Exploring node1: C
Exploring node1: D
Exploring node1: E
BFS traversal order: ['A', 'B', 'C', 'D', 'E']
\end{verbatim}
\end{tcolorbox}
\vspace{1em}

\section{Experiment 2: Depth First Search (DFS)}
\textbf{Source File:} \texttt{DFS.py} \\[4pt]
\textbf{Objective:} To implement Depth-First Search (DFS) graph traversal using a stack to explore each branch as deeply as possible before backtracking. \\[4pt]
\textbf{Input Parameters:} 4 edges: (A B), (A C), (B D), (C E); Starting node: A \\[6pt]
\begin{tcolorbox}[colback=gray!5!white,colframe=blue!70!black,title={\textbf{Program Output: Depth First Search (DFS)}}]
\begin{verbatim}
---Build Your Graph---
How many edges (connection) does your graph has ?Enter each edge separated by a spacce(e.g., A B):
Edge 1Edge 2Edge 3Edge 4Enter the starting node for DFS: 
Your Graph Dictionary: {'A': ['B', 'C'], 'B': ['A', 'D'], 'C': ['A', 'E'], 'D': ['B'], 'E': ['C']}
Staring DFS Traversal...
Exploring node1: A
Exploring node1: C
Exploring node1: E
Exploring node1: B
Exploring node1: D
DFS traversal order: ['A', 'C', 'E', 'B', 'D']
\end{verbatim}
\end{tcolorbox}
\vspace{1em}

\section{Experiment 3: Uniform Cost Search (UCS)}
\textbf{Source File:} \texttt{ucs.py} \\[4pt]
\textbf{Objective:} To implement Uniform Cost Search (UCS) using a priority queue to find the path with minimum cumulative edge cost from start to goal. \\[4pt]
\textbf{Input Parameters:} 5 nodes, 5 edges: (0 1: 2), (0 2: 4), (1 3: 7), (2 4: 3), (3 4: 1); Start: 0, Goal: 4 \\[6pt]
\begin{tcolorbox}[colback=gray!5!white,colframe=blue!70!black,title={\textbf{Program Output: Uniform Cost Search (UCS)}}]
\begin{verbatim}
===== UNIFORM COST SEARCH =====
Enter number of nodes: Enter number of edges: Enter edges as:
source destination cost
Enter starting node: Enter goal node: 
Goal Found!
Optimal Path: 0 -> 2 -> 4
Minimum Cost: 7
\end{verbatim}
\end{tcolorbox}
\vspace{1em}

\section{Experiment 4: Depth Limited Search (DLS)}
\textbf{Source File:} \texttt{dls.py} \\[4pt]
\textbf{Objective:} To implement Depth Limited Search (DLS) to mitigate the infinite loop vulnerability of DFS by constraining search depth. \\[4pt]
\textbf{Input Parameters:} 5 nodes, 4 edges: (0 1), (0 2), (1 3), (2 4); Start: 0, Goal: 4, Depth Limit: 3 \\[6pt]
\begin{tcolorbox}[colback=gray!5!white,colframe=blue!70!black,title={\textbf{Program Output: Depth Limited Search (DLS)}}]
\begin{verbatim}
===== DEPTH LIMITED SEARCH =====
Enter number of nodes: Enter number of edges: Enter edges (source destination):
Enter starting node: Enter goal node: Enter depth limit: 
Goal Found!
Path: 0 -> 2 -> 4
Depth: 2
\end{verbatim}
\end{tcolorbox}
\vspace{1em}

\section{Experiment 5: A* Search Algorithm}
\textbf{Source File:} \texttt{A_star.py} \\[4pt]
\textbf{Objective:} To implement A* Search Algorithm using evaluation function $f(n) = g(n) + h(n)$ combining exact cost and heuristic estimate. \\[4pt]
\textbf{Input Parameters:} 4 nodes (A, B, C, D) with heuristics h(A)=10, h(B)=8, h(C)=5, h(D)=0; 4 directed edges: A->B(1), A->C(4), B->D(5), C->D(2); Start: A, Goal: D \\[6pt]
\begin{tcolorbox}[colback=gray!5!white,colframe=blue!70!black,title={\textbf{Program Output: A* Search Algorithm}}]
\begin{verbatim}
=== A* Algorithm Input Setup ===
Enter the total number of nodes: 
Enter the heuristic value h(n) for each node:
Node name:  Heuristic h(A): Node name:  Heuristic h(B): Node name:  Heuristic h(C): Node name:  Heuristic h(D): 
Enter total number of directed edges:
Enter the edges in the format (from_node to_node weight)':
 Edge 1:  Edge 2:  Edge 3:  Edge 4: 
--- Path Finding ---
Enter the start node: Enter the goal node: 
== Result ===
Shortest path: A->B->D
Total Path Cost: 6.0
\end{verbatim}
\end{tcolorbox}
\vspace{1em}

\section{Experiment 6: Hill Climbing Search}
\textbf{Source File:} \texttt{HILL_CLIMBING.py} \\[4pt]
\textbf{Objective:} To implement the Hill Climbing local search algorithm for finding the local or global maximum of a mathematical objective function. \\[4pt]
\textbf{Input Parameters:} Starting value: 10, Step size: 0.1, Maximum iterations: 100 for objective function $f(x) = -(x^2) + 10$ \\[6pt]
\begin{tcolorbox}[colback=gray!5!white,colframe=blue!70!black,title={\textbf{Program Output: Hill Climbing Search}}]
\begin{verbatim}
Enter the starting value: Enter the step size: Enter maximum iterations: 
Best position found:  1.8790524691780774e-14
Maximum Value =  10.0
\end{verbatim}
\end{tcolorbox}
\vspace{1em}

\section{Experiment 7: Travelling Salesman Problem (TSP)}
\textbf{Source File:} \texttt{TSP.py} \\[4pt]
\textbf{Objective:} To solve the Travelling Salesman Problem by computing the minimum cost Hamiltonian path visiting every city once. \\[4pt]
\textbf{Input Parameters:} 4 cities with cost matrix: [[0, 10, 15, 20], [10, 0, 35, 25], [15, 35, 0, 30], [20, 25, 30, 0]], Starting city: 0 \\[6pt]
\begin{tcolorbox}[colback=gray!5!white,colframe=blue!70!black,title={\textbf{Program Output: Travelling Salesman Problem (TSP)}}]
\begin{verbatim}
Enter the number of cities: Enter the cost matrix:
Enter the starting city (0 to 3): 
Optimal path: 0->1->2->3
Minimum cost: 45
\end{verbatim}
\end{tcolorbox}
\vspace{1em}

\section{Experiment 8: 8-Queens Problem}
\textbf{Source File:} \texttt{Eight_Queen.py} \\[4pt]
\textbf{Objective:} To solve the 8-Queens constraint satisfaction problem using Backtracking search on an $8 \times 8$ board. \\[4pt]
\textbf{Input Parameters:} None (Autonomous backtracking search placement on an 8x8 chessboard) \\[6pt]
\begin{tcolorbox}[colback=gray!5!white,colframe=blue!70!black,title={\textbf{Program Output: 8-Queens Problem}}]
\begin{verbatim}
Solution for 8-Queens Problem:

Q . . . . . . .
. . . . Q . . .
. . . . . . . Q
. . . . . Q . .
. . Q . . . . .
. . . . . . Q .
. Q . . . . . .
. . . Q . . . .
\end{verbatim}
\end{tcolorbox}
\vspace{1em}

\section{Experiment 9: Water Jug Problem}
\textbf{Source File:} \texttt{Water_jug.py} \\[4pt]
\textbf{Objective:} To solve the Water Jug problem using State-Space Search (BFS) to measure a target quantity of water using two jugs. \\[4pt]
\textbf{Input Parameters:} None (Jug 1 Capacity: 4L, Jug 2 Capacity: 3L, Target state: (4, 2)) \\[6pt]
\begin{tcolorbox}[colback=gray!5!white,colframe=blue!70!black,title={\textbf{Program Output: Water Jug Problem}}]
\begin{verbatim}
steps to reach the target state (4, 2):
(0, 0)
(0, 3)
(3, 0)
(3, 3)
(4, 2)
\end{verbatim}
\end{tcolorbox}
\vspace{1em}

\section{Experiment 10: Tic-Tac-Toe Game}
\textbf{Source File:} \texttt{Tic_Tac_Toe.py} \\[4pt]
\textbf{Objective:} To implement a two-player interactive Tic-Tac-Toe game board with state representation and win/draw validation. \\[4pt]
\textbf{Input Parameters:} Sequential row and column moves: (0,0), (1,0), (0,1), (1,1), (0,2) yielding a win for Player X \\[6pt]
\begin{tcolorbox}[colback=gray!5!white,colframe=blue!70!black,title={\textbf{Program Output: Tic-Tac-Toe Game}}]
\begin{verbatim}
Player X's turn:
  |   |  
---------
  |   |  
---------
  |   |  
---------
Enter row (0-2): Enter column (0-2): 
Player O's turn:
X |   |  
---------
  |   |  
---------
  |   |  
---------
Enter row (0-2): Enter column (0-2): 
Player X's turn:
X |   |  
---------
O |   |  
---------
  |   |  
---------
Enter row (0-2): Enter column (0-2): 
Player O's turn:
X | X |  
---------
O |   |  
---------
  |   |  
---------
Enter row (0-2): Enter column (0-2): 
Player X's turn:
X | X |  
---------
O | O |  
---------
  |   |  
---------
Enter row (0-2): Enter column (0-2): X | X | X
---------
O | O |  
---------
  |   |  
---------

Player X wins!
\end{verbatim}
\end{tcolorbox}
\vspace{1em}

\section{Experiment 11: Minimax Algorithm}
\textbf{Source File:} \texttt{minmax.py} \\[4pt]
\textbf{Objective:} To implement the Minimax decision rule for adversarial game tree evaluation. \\[4pt]
\textbf{Input Parameters:} 8 leaf node values: 3 5 6 9 1 2 0 -1 at tree height 3 \\[6pt]
\begin{tcolorbox}[colback=gray!5!white,colframe=blue!70!black,title={\textbf{Program Output: Minimax Algorithm}}]
\begin{verbatim}
Enter 8 leaf node values: The optimal value is: 5
\end{verbatim}
\end{tcolorbox}
\vspace{1em}

\section{Experiment 12: Alpha-Beta Pruning}
\textbf{Source File:} \texttt{Alpha_Beta_pruning.py} \\[4pt]
\textbf{Objective:} To implement Alpha-Beta Pruning to prune non-promising branches from the Minimax game tree and reduce search complexity. \\[4pt]
\textbf{Input Parameters:} 8 leaf node values: 3 5 6 9 1 2 0 -1 with alpha=-infinity, beta=+infinity \\[6pt]
\begin{tcolorbox}[colback=gray!5!white,colframe=blue!70!black,title={\textbf{Program Output: Alpha-Beta Pruning}}]
\begin{verbatim}
Enter 8 leaf node values: 
Optimal value: 5
\end{verbatim}
\end{tcolorbox}
\vspace{1em}

\section{Experiment 13: Blocks World Problem}
\textbf{Source File:} \texttt{block_world.py} \\[4pt]
\textbf{Objective:} To solve the Blocks World planning problem using state-space BFS search to transition from initial block stacks to goal configuration. \\[4pt]
\textbf{Input Parameters:} 3 blocks (A, B, C); Initial: A on Table, B on A, C on Table; Goal: A on Table, B on Table, C on B \\[6pt]
\begin{tcolorbox}[colback=gray!5!white,colframe=blue!70!black,title={\textbf{Program Output: Blocks World Problem}}]
\begin{verbatim}
===== BLOCK WORLD PROBLEM =====
Enter number of blocks: 
Enter block names:
Block 1: Block 2: Block 3: 
===== INITIAL STATE =====
Enter the position of each block.
Enter 'Table' if the block is on the table.
Where is A initially? Where is B initially? Where is C initially? 
===== GOAL STATE =====
Enter the desired position of each block.
Where should A be in the goal state? Where should B be in the goal state? Where should C be in the goal state? 
Initial State:
B -> A
C -> Table
A -> Table

Goal State:
B -> Table
C -> B
A -> Table

===== SOLUTION =====
Step 1: Move B from A to Table
Step 2: Move C from Table to B

Total moves: 2
\end{verbatim}
\end{tcolorbox}
\vspace{1em}

\end{document}
